\documentclass[a4paper,11pt]{article}
\usepackage{fontspec}
\setmainfont{cmun}[Path=app/backend/assets/fonts/cmu/,Extension=.otf,UprightFont=*rm,BoldFont=*bx,ItalicFont=*ti,BoldItalicFont=*bi]
\setsansfont{cmun}[Path=app/backend/assets/fonts/cmu/,Extension=.otf,UprightFont=*rm,BoldFont=*bx,ItalicFont=*ti,BoldItalicFont=*bi]
\usepackage[russian]{babel}
\usepackage{amsmath,amssymb,graphicx,booktabs,tabularx,longtable,array}
\usepackage[margin=22mm,headheight=15pt]{geometry}
\usepackage[table]{xcolor}
\usepackage{enumitem,fancyhdr,titlesec,hyperref}
\definecolor{accent}{HTML}{7050C9}
\definecolor{ink}{HTML}{202638}
\definecolor{soft}{HTML}{F3F0FA}
\hypersetup{colorlinks=true,urlcolor=accent,linkcolor=accent,pdftitle={ЛистОК — документация проекта},pdfauthor={Команда ЛатехТим}}
\setlength{\parindent}{0pt}
\setlength{\parskip}{6pt}
\setlist{itemsep=4pt,topsep=4pt,leftmargin=5mm}
\renewcommand{\arraystretch}{1.28}
\pagestyle{fancy}\fancyhf{}
\fancyhead[L]{\small\sffamily\color{accent}ЛистОК}
\fancyhead[R]{\small\sffamily ЛатехТим · Идея фикс 2026}
\fancyfoot[L]{\small Документация проекта · 7 ноября 2026}
\fancyfoot[R]{\thepage}
\renewcommand{\headrulewidth}{0.3pt}
\titleformat{\section}{\Large\sffamily\bfseries\color{ink}}{\thesection}{0.7em}{}
\titleformat{\subsection}{\large\sffamily\bfseries\color{ink}}{}{0pt}{}
\newcommand{\note}[1]{\par\smallskip\colorbox{soft}{\parbox{\dimexpr\linewidth-2\fboxsep}{\small #1}}\par\smallskip}
\newcommand{\shot}[2]{\begin{center}\includegraphics[width=#1\linewidth]{#2}\end{center}}
\begin{document}
\thispagestyle{empty}
{\sffamily\color{accent}ХАКАТОН «ИДЕЯ ФИКС» · АРСЕНАЛ МАСТЕРА}
\vspace{18mm}

{\sffamily\bfseries\fontsize{42}{48}\selectfont ЛистОК}

{\sffamily\LARGE Рабочий лист, урок, результат}
\vspace{6mm}

\large Сервис для учителя: подготовка заданий, печать, проверка работ
и статистика класса в одном кабинете.\normalsize

\vspace{10mm}
\textbf{Команда «ЛатехТим»}

Капитан: Николай Семченков.

Участники: Максим Скоков, Дмитрий Мазка, Таисия Алексинская,
Анастасия Почкина.

Сборка и публикация: GitHub \texttt{pavelneumoin}.

\vspace{8mm}
\shot{0.95}{docs/images/teacher-result.png}

\vfill
\href{https://max.ru/ideya_fiks_126_bot}{Бот MAX: ideya\_fiks\_126\_bot}\hfill
\href{https://github.com/pavelneumoin/worksheet-creator}{Репозиторий проекта}

\small Документация действующего прототипа. Версия от 7 ноября 2026 года.
Демонстрационные данные вымышлены.\normalsize
\newpage
\section{Проблема, цель и задачи}
\subsection{Почему мы делаем этот проект}
Подготовка обычного урока распадается на много действий: найти задания,
перепечатать формулы, сверстать лист, подготовить ответы, раздать материал,
проверить работы и перенести результаты в таблицу. Даже удачный материал
сложно повторно использовать, если его версия и результаты класса хранятся отдельно.

\textbf{Цель проекта} — связать подготовку рабочего листа, выдачу классу
и проверку результатов в понятный учителю процесс. Основная функция —
создание качественного печатного материала; проверка и аналитика продолжают этот процесс.

\subsection{Задачи}
\begin{enumerate}
\item Принимать фотографии, PDF и текст, сохранять материал в личной библиотеке.
\item Формировать аккуратные математические листы LaTeX с CMU, ответами
и вариантами разной сложности.
\item Поддерживать классы, постоянные коды учеников, даты и неизменяемые
версии выданных работ.
\item Принимать пачки сканов и личные отправления учеников; выделять сомнения
и оставлять подтверждение оценки учителю.
\item Показывать результаты по заданиям, темам, классам и школе, выгружать отчёты.
\item Дать проверяющему готовый демонстрационный сценарий и воспроизводимый запуск.
\end{enumerate}

\subsection{Кому предназначен ЛистОК}
\begin{tabularx}{\linewidth}{@{}p{25mm}X@{}}
\toprule
\textbf{Роль}&\textbf{Основная задача}\\\midrule
Учитель&Подготовить, выдать и проверить работу, увидеть затруднения класса.\\
Ученик&Открыть свою работу, отправить решение, увидеть замечание и исправить его.\\
Завуч&Посмотреть охват проверки и результаты подключённых учителей по темам.\\
Проверяющий&Пройти готовый учебный сценарий, открыть файлы и проверить поведение.\\
\bottomrule
\end{tabularx}

\note{Рабочая гипотеза: связанный процесс уменьшает повторный ввод и переключение
между инструментами. Экономию времени и точность распознавания ещё нужно измерить
на реальных уроках; в проекте нет выдуманных показателей эффективности.}
\newpage
\section{Сценарий учителя}
\subsection{Подготовка и печать}
На первом экране учитель добавляет материал и выбирает оформление. Задание
обрабатывается в фоне; кабинет остаётся доступен. Результат сохраняется в личной
библиотеке. Публикация для коллег выполняется отдельным действием.

Вверху готового результата находятся назначение классу и файлы: рабочий лист,
ответы, комплект учителя и бланк A4. Варианты «Другой», «Проще» и «Сложнее»
расположены рядом. Длинный просмотр PDF и исходники свёрнуты ниже.

\shot{0.96}{docs/images/teacher-result.png}

\subsection{Выдача и проверка}
Учитель выбирает класс и дату. Назначение хранит конкретную версию листа
и состав класса на момент выдачи. Изменение материала в библиотеке не меняет
задание, которое уже получили ученики.

После урока доступны два пути: общий PDF из МФУ или индивидуальная отправка
ученика. В пачке можно последовательно проверять только сомнительные ответы
разных учеников, увеличивать фрагмент и ставить частичные баллы. Перед включением
в статистику результат каждой работы подтверждается.

\note{Главные действия доступны с телефона. Проверены ширины 360, 390 и 430 px.
Это проверка адаптивной веб-страницы в Edge, а не испытание встроенного MAX
на настоящих Android- и iOS-устройствах.}
\newpage
\section{Печатные материалы и математика}
\subsection{Один материал — несколько файлов}
\begin{tabularx}{\linewidth}{@{}p{44mm}X@{}}
\toprule\textbf{Файл}&\textbf{Назначение}\\\midrule
Рабочий лист&Условия и место для решения; без ключа ответов.\\
Ответы&Отдельный учительский файл с ответами и пояснениями.\\
Комплект учителя&Лист и ответы в одном PDF.\\
Бланк ответов&Номера заданий и клетки на одной странице A4.\\
Персональные бланки&Код ученика, версия и QR назначения для пакетного сканирования.\\
\bottomrule\end{tabularx}

\subsection{Бланк для быстрой проверки}
Если на многостраничном листе нужны только итоговые ответы, ученик переносит их
на отдельную страницу. Учитель сканирует эти страницы. Для задач с доказательством
или развёрнутым вычислением сохраняется загрузка полного решения.
До назначения доступен общий бланк; после назначения — персональный комплект.
Короткий код и QR служат сопоставлению работ и не являются паролем ученика.

\subsection{Типографика}
Документы собираются XeLaTeX. Текст — Computer Modern Unicode, математические
выражения — LaTeX. Проверены дроби, корни, степени, логарифмы, тригонометрия,
системы и русские условия. Например:
\[
2\sin^2 x-3\sin x+1=0,\qquad
\log_2(x-1)+\log_2(x-3)=3.
\]
У второго уравнения область допустимых значений $x>3$, поэтому из корней
преобразованного квадратного уравнения $-1$ и $5$ остаётся только $5$.

\subsection{Демонстрационные материалы}
В комплект входят тригонометрические уравнения с отбором корней, логарифмы
с ОДЗ и длинные задачи по физике: движение с трением, нагрев с потерями энергии,
электрическая цепь. Сложная физика занимает две страницы, ответы — одну.
Готовые PDF проверяются отдельно от ответа модели: наличие файла само по себе
не подтверждает правильность математики.

\note{Учитель сверяет условия и ключ перед выдачей. Неясный текст фотографии
нельзя молча заменять догадкой. Введённое учителем неполное условие может получить
предложение уточнения, которое отмечается для проверки.}
\newpage
\section{Кабинет ученика и исправления}
\begin{minipage}[t]{0.61\linewidth}
Учитель выдаёт персональную ссылку из списка класса. Ученик входит через MAX
и принимает приглашение. Ссылка действует 24 часа, используется один раз;
новая ссылка отменяет старую.

В кабинете доступны только собственные задания, PDF без ключей, подтверждённые
баллы и комментарии. Можно раскрыть отдельное задание и посмотреть свой результат.

\subsection{Повторная сдача}
\begin{enumerate}
\item Учитель возвращает работу с комментарием.
\item Ученик прикрепляет исправление: PDF либо фото страниц.
\item Новая попытка появляется у учителя в «От учеников» и в журнале класса.
\item После проверки новый результат подтверждается; история остаётся.
\end{enumerate}
До этого момента прежний подтверждённый балл продолжает участвовать в статистике.
Если результат изменён другим сеансом, требуется сверить актуальную версию.

Принимаются один PDF до десяти страниц или до десяти фотографий: не более
15 МБ на файл и 20 МБ на отправку. Число попыток ограничено десятью.
\end{minipage}\hfill
\begin{minipage}[t]{0.34\linewidth}
\vspace{0pt}\includegraphics[width=\linewidth]{docs/images/student-mobile.png}
\end{minipage}

\subsection{Защита доступа}
Привязка проверяется сервером для каждого запроса. Ученик не получает список
класса, чужие сканы, ключи и учительские экспорты. После отзыва привязки старые
ссылки на файлы не дают доступа. Демо и настоящие классы изолированы друг от друга.

\note{Правки проверки сохраняются автоматически в черновик. При возврате работы
бот уведомляет ученика в MAX. Подтверждённая оценка меняется только по решению учителя.
Создать новый лист можно голосовым сообщением: сначала подтвердите распознанный запрос.}
\newpage
\section{Проверка решений и статистика}
\subsection{Разделение распознавания и оценки}
На первом этапе модель переписывает видимое решение без ключа ответов. На втором
получает условие, ключ и расшифровку: проверяет ограничения, обоснование,
вычисления и предлагает балл с объяснением. Программа требует основание из
распознанного текста. Если текст не прочитан или основание отсутствует,
предложенный балл убирается. Сам анализ оценку не меняет.

Для развёрнутой математики используются краткие авторские ориентиры по темам.
Учтены общие принципы материалов ФИПИ: полнота обоснования, допустимость
корректных альтернативных решений, различение метода и вычислительной ошибки.
Критерии экзамена не переносятся автоматически на произвольный рабочий лист.

\subsection{Статистика класса и школы}
Процент выполнения рассчитывается только по подтверждённым результатам:
\[
P=100\,\frac{\sum\text{набранных баллов}}{\sum\text{максимальных баллов}}.
\]
Отдельно показан охват — доля подтверждённых работ среди назначенных.
Несданные и непроверенные работы не считаются нулевыми.
Завуч видит данные учителей, которые присоединились к школе с согласием на сводку.
Есть фильтры по периоду, классу, теме и разделу; PDF и CSV содержат подробные таблицы.

\shot{0.94}{docs/images/school-analytics.png}
\newpage
\section{Устройство и хранение}
\begin{tabularx}{\linewidth}{@{}p{39mm}X@{}}
\toprule\textbf{Часть}&\textbf{Ответственность}\\\midrule
Веб-кабинет&HTML, CSS, JavaScript; интерфейсы учителя, ученика и завуча.\\
Flask API&Вход, права доступа, материалы, классы, работы, отчёты.\\
Очередь в SQLite&Постоянное состояние заданий, повторы и защита от дублей.\\
Обработчик модели&Распознавание, варианты, рекомендации и анализ решений.\\
LaTeX-обработчик&Печать в отдельном контейнере без сети и ключей модели.\\
MAX&Webhook бота, вход пользователя, быстрые действия в чате.\\
Docker Compose&Воспроизводимый запуск сервисов и томов данных.\\
\bottomrule\end{tabularx}

\subsection{Данные и версии}
Материал принадлежит учителю. Сохранённая печатная версия связана с назначением;
назначение содержит снимок состава класса. Проверка и история попыток ссылаются
на это назначение. Версии записей защищают от случайной перезаписи при работе
на нескольких устройствах. Повтор запроса отправки не создаёт дубликат.

Имена из импорта списка выдаются учителю в локальной таблице соответствия,
в базе класса используются коды. Сами сканы могут содержать написанные учеником
имена, поэтому доступ к ним ограничен. Сканы хранятся семь дней; результаты
и комментарии остаются. Это техническое правило прототипа, а не заключение
о выполнении всех требований к школьной информационной системе.

\subsection{Сбой и восстановление}
Длительная обработка не выполняется внутри веб-запроса. Результат сохраняется
в очереди и библиотеке. При ошибке задача показывает причину и допускает
контролируемый повтор. Сохранённые части не должны генерироваться заново без нужды.

Резервная копия содержит базу, связанные PDF и ещё не истёкшие сканы.
Восстановление проверяет контрольные суммы, целостность SQLite и ссылки на файлы.
Активные сессии и приглашения отзываются. Секреты хранятся отдельно от кода и образа.

\subsection{Встроенный MAX}
Вход через конкурсного бота работает. Для запуска страницы как нативного
мини-приложения владелец бота должен привязать HTTPS-адрес. У команды есть токен,
а управление конкурсным ботом остаётся у организаторов. Это внешний шаг запуска.
\newpage
\section{Проверка и соответствие конкурсу}
\subsection{Что проверяется}
\begin{itemize}
\item 287 Python-тестов и 26 Node-тестов: доступ, версии, очередь, ответы,
приглашения, исправления, экспорты, повторные запросы и резервные копии.
\item Настоящая компиляция PDF: CMU, сложная математика, бланки на 30 учеников,
таблицы и отчёты. Визуальная проверка страниц после рендеринга.
\item Адаптивная страница, доставка учебной попытки, сохранение старого балла,
переходы между учительским и ученическим сценариями.
\item Отдельный запрос к действующему API на печатном учебном примере:
полное тригонометрическое решение и пропущенная ОДЗ в логарифмическом.
\item CI: сборка контейнера, тесты без сети, запуск сервисов, перезапуск,
проверка готовности, экспорт образа и контрольной суммы.
\end{itemize}

\note{Результат на печатном учебном примере не измеряет качество работы
с почерком. Для такой оценки нужна отдельная обезличенная выборка и разметка учителей.}

\subsection{Критерии полуфинала}
Основание: положение «Идея фикс», редакция 06.10.2026, раздел 10.4,
страницы 8–10. Четыре критерия по пять баллов; баллы выставляет жюри.

\begin{tabularx}{\linewidth}{@{}p{43mm}X@{}}
\toprule\textbf{Критерий}&\textbf{Проверяемое свидетельство}\\\midrule
Соответствие требованиям&Работающий бот, API и веб-кабинет; путь от материала
до выдачи и отчёта. Привязка встроенного MAX отмечена как внешний шаг.\\
Пользовательский опыт&Файлы сразу в результате, короткие действия,
мобильная проверка и контекстные команды бота.\\
Стабильность&Очередь, ошибки, идемпотентность, ограничения файлов,
проверки доступа и восстановление копии.\\
Мобильность архитектуры&Отдельные модули, адаптивный интерфейс,
MAX, API модели, LaTeX, переносимый Docker-запуск.\\
\bottomrule\end{tabularx}

Срок из последнего сообщения организаторов: 8 ноября, 23:59 МСК.
Работу подаёт капитан отдельно. Публикация кода и образа не означает подачи заявки.
\newpage
\section{Инструкция проверяющему и развитие}
\subsection{Показ за несколько минут}
\begin{enumerate}
\item Открыть бота MAX, отправить \texttt{/start}, войти в кабинет и выбрать демо.
\item Открыть тригонометрический лист, ответы и бланк A4.
\item Посмотреть класс, назначение, пачку и один сомнительный ответ.
\item Перейти к школьной статистике: темы, разделы, охват и экспорт.
\item Открыть «От учеников» и «Кабинет ученика»: замечание, попытка, исправление.
\end{enumerate}
В новом демо подготовлены шесть классов, в текущих списках по 30 учеников;
старые назначения сохраняют исходный состав. Материалы включают логарифмы,
тригонометрию и текстовую физику. Данные явно учебные и доступны без загрузки своих файлов.

\subsection{Запуск из исходников}
Установить Docker, скопировать \texttt{.env.example} в \texttt{.env} и заполнить
параметры из \texttt{RUN.md}. Команды из корня репозитория:
\begin{verbatim}
docker build -t listok:night7 .
docker compose up -d --wait
curl --fail http://127.0.0.1:3005/api/health/ready
\end{verbatim}
В GitHub Releases лежат готовый архив образа и SHA256 после успешного CI.
Для публичного запуска используются дополнительные compose-файлы HTTPS и бота.
Версии библиотек закреплены в \texttt{requirements.txt}.

\subsection{Следующие этапы}
\begin{enumerate}
\item Подключить встроенный MAX и проверить Android/iOS: возврат, камера,
скачивание PDF и восстановление после сворачивания.
\item Собрать набор обезличенных рукописных работ с эталонной проверкой.
Измерить ошибки отдельно для распознавания и оценки по темам.
\item Улучшить выравнивание бланков и выделение фрагмента решения; добавить
критерии конкретного задания и разный максимальный балл.
\item Проверить голосовое создание листов на записях нескольких учителей;
измерить ошибки терминов, чисел и общее время подготовки.
\end{enumerate}
\small
Источники методических ориентиров:
\href{https://doc.fipi.ru/oge/dlya-predmetnyh-komissiy-subektov-rf/2026/mr_oge_matematika_2026.pdf}{ФИПИ, ОГЭ, математика, 2026};
\href{https://doc.fipi.ru/ege/dlya-predmetnyh-komissiy-subektov-rf/2026/matematika_mr_ege_2026.pdf}{ФИПИ, ЕГЭ, математика, 2026}.
Использован краткий авторский пересказ общих принципов.
\end{document}
